\chapter{绪论}

\section{课题研究的背景及意义}

高校信息化建设正在从单一业务系统建设转向平台化、智能化和协同化发展。校园网作为各类教学平台、办公系统、资源库和安全管理系统的基础承载环境，其稳定性与安全性会直接影响学校日常运行。随着移动终端、在线课程、实验平台和数据服务的持续增长，传统扁平化网络结构已经难以满足现代校园的业务需求。

本课题以高校校园网建设为应用场景，围绕网络拓扑、地址规划、路由策略和安全访问控制等内容开展设计。通过将课程中学习的数据通信理论与工程实践结合，可以训练学生综合运用网络规划、设备配置、系统分析和技术文档编写等能力。

\subsection{课题背景}

现代校园网络通常需要同时承载教学、科研、办公、宿舍和公共服务等多类业务。不同区域在带宽、访问权限、网络安全和运维管理方面存在差异，因此需要在整体架构上进行统一规划，并在局部接入层面保留灵活配置能力。

\subsection{项目的意义}

合理的校园网建设方案能够提高网络资源利用率，降低故障扩散风险，并提升后期运维效率。对于应用型本科毕业设计而言，该类课题既能体现真实工程约束，也便于通过拓扑图、配置说明、测试表和安全策略表展示设计成果。

\section{与本课题相关的国内外现状}

当前校园网络建设普遍采用分层架构，并通过虚拟局域网、链路冗余、动态路由和集中认证等技术实现业务隔离与统一管理。随着智慧校园建设推进，网络系统还需要兼顾物联网设备接入、无线覆盖优化和数据安全治理等新要求。

\section{论文组织}

本文共分为五章。第一章介绍课题背景、研究意义和论文组织；第二章进行系统分析与网络方案设计；第三章说明关键功能与安全策略；第四章给出系统测试方案；第五章进行总结与展望。
